%% =================================================================
%% Wei, Mei, Zhao, Xiao, Sun, Zhang, Guo.
%% "Nondestructively identifying the mechanical behavior of soft
%%  tissues using surface deformation with an explicit inverse
%%  approach."
%% Appl. Math. Modelling 134 (2024) 126--147.
%%
%% Reconstruction of page 3 (printed p. 128) as study notes.
%% Prose: extracted verbatim from the PDF text layer via pdftotext
%%        (soft hyphens and line breaks cleaned up only).
%% Math:  hand-transcribed from the rendered page image — Eqs. (5)–(7)
%%        including the Cox–de Boor nested-cases recursion in Eq. (6)
%%        and the closed B-spline control-point definition in Eq. (7).
%% Figure 1: placeholder box with caption + visual description.
%% =================================================================

\documentclass[11pt]{article}

\usepackage[margin=1in]{geometry}
\usepackage{amsmath, amssymb}
\usepackage{mathrsfs}
\usepackage{fancyhdr}

\pagestyle{fancy}
\fancyhf{}
\lhead{\small Y.~Mei et al.}
\rhead{\small Applied Mathematical Modelling 134 (2024) 126--147}
\cfoot{\small \thepage}
\renewcommand{\headrulewidth}{0.4pt}

\setcounter{page}{128}
\setcounter{equation}{4}  % next equation is (5)

\begin{document}

\noindent
where the superscript $h$ denotes the finite element subspace.
$\mathscr{M}^h$, $\ell^h$ and $\mathscr{P}^h$ are the finite-element
subspaces of the function space
$\mathscr{M} \equiv \{\mathbf{u} \mid u_i \in H^1(\Omega);\,
u_i = g_i \text{ on } \Gamma_g\}$,
$\ell \equiv \{\mathbf{w} \mid w_i \in H^1(\Omega);\,
w_i = 0 \text{ on } \Gamma_g\}$ and
$\mathscr{P} \subseteq L^2(\Omega)$, respectively. The third term is
the stabilized term to address the volumetric locking induced by the
incompressibility where $NE$ is the total number of elements in the
discretized domain. $\tau = \dfrac{\alpha h^{2}}{2\mu}$ represents the
stabilization factor, where $h$ is the characteristic element length.
$\alpha$ is set to $1.0$ [30].

\subsection*{\textit{Topological description of nonhomogeneous elastic property distribution}}

\noindent
To describe the topology of the nonhomogeneous elastic property
distribution, we introduce the B-spline function to express a curve:
\begin{equation}
\mathbf{C}(t) = \sum_{i=0}^{n+1} N_{i,p}(t)\,\mathbf{P}_i,
\qquad (a \leq t \leq b)
\end{equation}
where $N_{i,p}(t)$ is the $p$-th order B-spline basis function
generated by a knot vector $\{t_0, t_1, \ldots, t_{n+2+p}\}$.
$\mathbf{P}_i\,(i = 0, \ldots, n+1)$ is the $i$-th control point.

The function $N_{i,p}(t)$ can be determined by Eq.~(6):
\begin{equation}
\begin{cases}
N_{i,0}(t) =
\begin{cases}
1, & \text{if } t_i \leq t \leq t_{i+1} \\
0, & \text{otherwise}
\end{cases}
& (i = 0, \ldots, n+1;\ p = 0) \\[3.5ex]
N_{i,p}(t) = \dfrac{t - t_i}{t_{i+p} - t_i}\,N_{i,p-1}(t)
+ \dfrac{t_{i+p+1} - t}{t_{i+p+1} - t_{i+1}}\,N_{i+1,p-1}(t),
& (i = 0, \ldots, n+1;\ p \geq 1)
\end{cases}
\end{equation}

To capture various types of nonhomogeneous elastic property
distributions, we employ B-splines to approximate both open and
closed curves, as depicted in Fig.~1. The open curve is utilized to
represent the interface shape between adjacent layers in a layered
composite, while the closed curve represents the shape of inclusions
embedded within a soft solid. To avoid the self-intersection of the
closed B-spline curve, we define the control points of the closed
B-spline curve in Eq.~(7):
\begin{equation}
\mathbf{P}_i =
\begin{cases}
\left(
x_o + d_i \sin\!\left(i\theta - \dfrac{\theta}{2}\right),\;
y_o + d_i \cos\!\left(i\theta - \dfrac{\theta}{2}\right)
\right),
& \left(\theta = \dfrac{2\pi}{n};\ i = 1, \ldots, n\right) \\[3ex]
\mathbf{P}_0 = \mathbf{P}_{n+1} = \dfrac{\mathbf{P}_1 + \mathbf{P}_n}{2}
\end{cases}
\end{equation}
where $(x_o, y_o)$ is the location of the center of the closed
B-spline curve, $d_i$ is the distance between the $i$-th control
point and the center of the closed curve.

With the curve expressed by the B-spline function, we can define
each nonhomogeneous region in the problem domain by

% ---------- Figure 1 placeholder ----------
\vspace{1em}
\begin{center}
\fbox{\parbox{0.92\textwidth}{\centering\vspace{1em}
\textbf{[Figure 1 — placeholder]}\\[0.8em]
\textbf{(a) Open curve}: eight control points
$\mathbf{P}_0, \mathbf{P}_1, \ldots, \mathbf{P}_7$ connected by dashed
blue segments, with a smooth orange B-spline curve passing through
the polygon.\\[0.4em]
\textbf{(b) Closed curve}: nine control points
$\mathbf{P}_0 = \mathbf{P}_9, \mathbf{P}_1, \ldots, \mathbf{P}_8$
distributed around a center $(x_o, y_o)$, with radial distances
$d_1, d_2, \ldots, d_8$ and angular spacing $\theta = 2\pi/n$.
The closed B-spline curve (orange) is inscribed in the dashed
control polygon, with $\mathbf{P}_0 = \mathbf{P}_9 = (\mathbf{P}_1 + \mathbf{P}_8)/2$.
\vspace{1em}}}
\end{center}
\vspace{0.4em}
\begin{center}
\textbf{Fig. 1.} Two examples of the curve expressed by B-spline functions.
\end{center}

\end{document}
